\documentclass[acmsmall,screen,nonacm]{acmart}

\usepackage{bcprules}
\usepackage{mathtools}
\usepackage{float}
\usepackage[nameinlink,noabbrev]{cleveref}

\settopmatter{printacmref=false,printccs=false,printfolios=true}
\raggedbottom

\newtheorem{theorem}{Theorem}
\newtheorem{lemma}[theorem]{Lemma}

\newcounter{RuleRef}
\newcommand{\ruledef}[1]{%
  \refstepcounter{RuleRef}\textsc{#1}\label{rule:#1}}
\newcommand{\ruleuse}[1]{%
  \hyperref[rule:#1]{\textsc{#1}}}
\newcommand{\ruleref}[1]{%
  \textsc{(\hyperref[rule:#1]{#1})}}

\newcommand{\lambdapcc}{\ensuremath{\lambda_{\mathsf p}^{\mathsf{cc}}}}
\newcommand{\ctxempty}{\varnothing}
\newcommand{\cempty}{\varnothing}
\newcommand{\cset}[1]{\{#1\}}
\newcommand{\capt}[2]{#1^{#2}}
\newcommand{\sing}[1]{#1.\mathsf{type}}
\newcommand{\pairshape}[4]{\langle #1\!:\!#2;\,#3\!:\!#4\rangle}
\newcommand{\pairval}[3]{\langle #1;\,#2=#3\rangle}
\newcommand{\subst}[3]{#1[#2/#3]}
\newcommand{\wf}{\mathrel{\mathsf{wf}}}
\newcommand{\uses}{\mathbin{\triangleright}}
\newcommand{\sstep}{\longrightarrow}
\newcommand{\steps}{\longrightarrow^{*}}
\newcommand{\resolve}{\Downarrow}
\newcommand{\loc}[1]{\mathsf{loc}(#1)}
\newcommand{\typeref}[1]{\mathsf{type}(#1)}
\newcommand{\capref}[1]{\mathsf{capture}(#1)}
\newcommand{\refdef}[1]{\mathsf{ref}(#1)}
\newcommand{\tcoerce}{\mathrel{\rightsquigarrow_{\mathsf t}}}
\newcommand{\scoerce}{\mathrel{\rightsquigarrow_{\mathsf s}}}
\newcommand{\coerce}{\mathrel{\rightsquigarrow}}
\newcommand{\extends}{\mathrel{\preccurlyeq_{\mathsf{alloc}}}}
\newcommand{\semleq}{\mathrel{\preceq}}
\newcommand{\qual}{\mathsf{qual}}
\newcommand{\valid}{\mathrel{\mathsf{valid}}}
\newcommand{\returns}{\mathrel{\mathsf{returns}}}
\newcommand{\plateheader}[2]{%
  \par\addvspace{1em}%
  \noindent\textbf{#1}%
  \if\relax\detokenize{#2}\relax\else
    \hfill{\setlength{\fboxsep}{3pt}\setlength{\fboxrule}{0.5pt}%
      \fbox{\(\strut #2\)}}%
  \fi\par
  \nobreak\addvspace{0.65em}%
}
\newcommand{\rulesetup}{%
  \small
  \smallrulenames
  \setlength{\afterruleskip}{1.45em plus 0.2em minus 0.1em}%
  \setlength{\labelminsep}{0.65em}%
  \renewcommand{\arraystretch}{1.15}%
  \renewcommand{\andalso}{\qquad}%
}

\title{Capture-Checked Paths through Dependent Pairs}
\author{Oliver Bra\v{c}evac}
\author{Martin Odersky}
\renewcommand{\shortauthors}{Bra\v{c}evac and Odersky}
\authorsaddresses{}

\begin{document}

\begin{abstract}
We extend a small calculus of path-dependent types based on immutable
dependent pairs with capture qualifiers, use sets, and abstract capture-set
members. Capture sets contain paths of arbitrary length. A pair member may
describe a term, a type interval, or a capture-set interval, and its signature
may depend on the pair's first component. Subcapturing expands a proper path
to the qualifier of its synthesized type, relates an abstract capture
selection to its declared bounds, and contracts proper projections and term
selections to their root.

The soundness proof interprets subtyping derivations as store-indexed
semantic coercions.  An annotated store records the qualifier assigned by
the introduction rule of each allocated value.  Qualifier-aware inhabitants
retain this introduction qualifier and a subcapturing derivation to the
qualifier of their assigned type.  Coercion action preserves the introduction
qualifier, including through dependent functions, pairs, and abstract
members.  The general pair rule is handled by applying its first-component
coercion at the stored component and instantiating its member coercion at the
same location.  Coercion action is well founded in the allocation order of
the immutable store.

The mechanization proves progress, preservation, and non-stuckness for
closed programs.  It also proves that both operands inspected by every
application are covered by the source use set, transported through preceding
allocations, and that the introduction qualifier of a returned value is
below the qualifier of its result type.
\end{abstract}

\keywords{path-dependent types, capture checking, dependent pairs,
subtyping, semantic coercions, type safety}

\maketitle

\section{Introduction}

Path-dependent types allow the type of a program phrase to mention a stable
path. Capture checking attaches a finite description of capabilities to a
value and tracks the capabilities used while evaluating a term. Both
mechanisms become more useful when paths may be longer than variables: a
component selected from a structure can carry its own type and its own
capture information.

The calculus studied here, \lambdapcc{}, combines these ideas in a small
call-by-value language of immutable dependent pairs. It retains the path,
singleton, abstract-type, function, and general pair-subtyping mechanisms of
the underlying path calculus. It adds three ingredients from capture
checking. First, every proper type is a shape qualified by a capture set.
Second, term typing produces a \emph{use set}, an upper bound on the
capabilities used to evaluate the term. Third, pairs may contain abstract
capture-set members with lower and upper bounds, in parallel with abstract
type members.

A qualified type is written \(\capt{S}{C}\). The shape \(S\) describes the
value, while \(C\) records capabilities retained by that value. A typing
judgment
\[
  \Gamma\vdash t:T\uses C
\]
assigns result type \(T\) and use set \(C\) to \(t\). Values have empty use
sets because their evaluation is complete. A function value is therefore
pure as an expression, while the qualifier on its function type records the
capabilities used by its body. This separation gives the usual distinction
between capabilities used now and capabilities retained for later use
\citep{xu2025box,boruch2023capture}.

Capture sets may contain a capability path \(p\) and an abstract capture
selection \(p.a\). We write the former as the singleton \(\cset p\); the
latter is itself a capture-set expression. The two forms have different
roles. If \(p\) synthesizes \(\capt{S}{C}\), subcapturing derives
\(\cset p<:C\). If the capture member \(p.a\) has bounds \(L..U\), then
\(L<:p.a<:U\). Proper projections and term selections contract to their
owner, for example \(\cset{p.a}<:\cset p\). There is no corresponding rule
\(p.a<:\cset p\): a capture member is a static definition whose contents
need not be owned by its receiver. This distinction prevents an arbitrary
concrete capture definition from being collapsed through the receiver's
qualifier.

The main subtype of interest is the general pair rule:
\begingroup
\rulesetup
\infrule[\ruleuse{ss-pair}]
  {\Gamma\vdash T<:T'
   \andalso
   \Gamma,x:T\vdash\mu<:\mu'}
  {\Gamma\vdash
    \pairshape{x}{T}{a}{\mu}<:
    \pairshape{x}{T'}{a}{\mu'}}
\endgroup
The member signature \(\mu\) can be a qualified type, a type interval, or a
capture interval. Its premise cannot be compiled to a closed semantic
coercion until a concrete pair supplies the location denoted by \(x\).
The proof retains this premise as a member closure. When a pair at location
\(z\) is observed, its runtime value exposes a first component \(y\) and a
member referent. The proof coerces \(y\), instantiates the member closure with
\(x\mapsto y\), and applies the resulting coercion to that referent.

The semantic interpretation retains both capture-related indices of source
typing.  A value carries its \emph{introduction qualifier}: the qualifier
fixed by the value rule before any enclosing subsumption.  For an abstraction
this is the free part of its body's use set; for a pair it is determined by
the stored components.  An introduction qualifier need not be minimal or
principal, since its source derivation may itself contain subsumption.
Runtime term evidence also retains the use set assigned to the current term.
These two indices support separate results for ordinary type safety, use
prediction, and capture prediction.

The rest of the paper presents the complete calculus, its operational
semantics, the coercion interpretation, and the soundness argument. The
presentation follows the formal judgments rather than the auxiliary
organization of the mechanization.

\section{The calculus}
\label{sec:calculus}

\subsection{Notation and syntactic categories}

Variables are written \(x,y,z\), and labels are written \(a,b\). Labels
inhabit one syntactic namespace. The kind of a selection is determined by
path synthesis, rather than by label typography. Paths \(p,q\) consist of variables, first
projections, and labelled selections. Only a path that synthesizes a
qualified type may occur as a term or as a singleton capture.

Capture sets are written \(C,D\), shapes \(R,S\), qualified types \(T,U\),
and generalized member signatures \(\mu,\nu\). A member signature is a
qualified term type, a shape interval, or a capture interval. We omit an
empty qualifier, writing \(S\) for \(\capt{S}{\cempty}\). The singleton
shape of a path is \(\sing p\), whereas \(\cset p\) is a singleton capture
set. The binder in a function scopes over its codomain. The binder in
\(\pairshape{x}{T}{a}{\mu}\) scopes over \(\mu\) and denotes the first
component. Capture-avoiding opening is written \(\subst{X}{p}{x}\) for any
type, shape, signature, or capture set \(X\).

Primes and numeric subscripts denote further metavariables of the same
category.  In rules that apply uniformly to type and capture intervals,
\(L\) and \(U\) denote the lower and upper bound; the member kind determines
whether they range over shapes or capture sets.

The term language is in monadic normal form. Applications take paths as
operands, and a pair stores a first-component variable together with a term
variable, shape definition, or capture-set definition. The runtime retains
shape and capture definitions so generalized path lookup has a uniform
referent.

\subsection{Static semantics}

The next three pages give the complete source syntax, static semantics,
runtime syntax, generalized path resolution, and CK machine.

\clearpage
\begingroup
\setlength{\abovedisplayskip}{5pt plus 1pt minus 1pt}
\setlength{\belowdisplayskip}{5pt plus 1pt minus 1pt}
\setlength{\intextsep}{5pt}
\captionsetup{skip=6pt}
\rulesetup
\setlength{\afterruleskip}{0.65em plus 0.1em minus 0.05em}
\renewcommand{\plateheader}[2]{%
  \par\addvspace{0.45em}%
  \noindent\textbf{#1}%
  \if\relax\detokenize{#2}\relax\else
    \hfill{\setlength{\fboxsep}{3pt}\setlength{\fboxrule}{0.5pt}%
      \fbox{\(\strut #2\)}}%
  \fi\par
  \nobreak\addvspace{0.45em}%
}

\begin{figure}[H]
\plateheader{Source syntax}{}
\noindent
\begingroup
\renewcommand{\arraystretch}{0.94}
\begin{minipage}[t]{0.48\linewidth}
\begin{tabular*}{\linewidth}{@{}r@{\ }l@{\extracolsep{\fill}}r@{}}
\(p,q ::= \) && \textbf{Path}\\
 & \(x\) & variable\\
 & \(p.1\) & first projection\\
 & \(p.a\) & member selection\\[0.35em]
\(C,D ::= \) && \textbf{Capture set}\\
 & \(\cempty\) & empty\\
 & \(C\cup D\) & union\\
 & \(\cset p\) & capability path\\
 & \(p.a\) & capture selection\\[0.35em]
\(T,U ::= \) && \textbf{Qualified type}\\
 & \(\capt{S}{C}\) & capturing\\[0.35em]
\(R,S ::= \) && \textbf{Shape}\\
 & \(\top\mid\bot\) & extrema\\
 & \((x:T)\to U\) & function\\
 & \(\pairshape{x}{T}{a}{\mu}\) & pair\\
 & \(\sing p\) & singleton\\
 & \(p.a\) & type selection
\end{tabular*}
\end{minipage}\hfill
\begin{minipage}[t]{0.48\linewidth}
\begin{tabular*}{\linewidth}{@{}r@{\ }l@{\extracolsep{\fill}}r@{}}
\(\mu,\nu ::= \) && \textbf{Member signature}\\
 & \(T\) & term member\\
 & \(S..R\) & type interval\\
 & \(C..D\) & capture interval\\[0.35em]
\(\delta ::= \) && \textbf{Member definition}\\
 & \(z\) & term definition\\
 & \(S\) & type definition\\
 & \(C\) & capture definition\\[0.35em]
\(t,u ::= \) && \textbf{Term}\\
 & \(p\) & path\\
 & \(\lambda(x:T).t\) & abstraction\\
 & \(\pairval{y}{a}{\delta}\) & pair\\
 & \(p\,q\) & application\\
 & \(\mathbf{let}\ x=t\ \mathbf{in}\ u\) & let\\[0.35em]
\(v ::= \) && \textbf{Value}\\
 & \(\lambda(x:T).t\mid\pairval{y}{a}{\delta}\) &\\[0.35em]
\(\Gamma ::= \) && \textbf{Context}\\
 & \(\ctxempty\mid\Gamma,x:T\) &
\end{tabular*}
\end{minipage}
\endgroup

\plateheader{Path synthesis}{\Gamma\vdash p\Rightarrow\mu}
\noindent
\begin{minipage}[t]{0.47\linewidth}
\infrule[\ruledef{p-var}]
  {\Gamma(x)=T}
  {\Gamma\vdash x\Rightarrow T}

\infrule[\ruledef{p-fst}]
  {\Gamma\vdash p\Rightarrow
    \capt{\pairshape{x}{T}{a}{\mu}}{C}}
  {\Gamma\vdash p.1\Rightarrow T}
\end{minipage}\hfill
\begin{minipage}[t]{0.51\linewidth}
\infrule[\ruledef{p-sel}]
  {\Gamma\vdash p\Rightarrow
    \capt{\pairshape{x}{T}{a}{\mu}}{C}}
  {\Gamma\vdash p.a\Rightarrow\subst{\mu}{p.1}{x}}

\infrule[\ruledef{p-miss}]
  {\Gamma\vdash p\Rightarrow
     \capt{\pairshape{x}{T}{b}{\nu}}{C}
   \\
   \Gamma\vdash p.1.a\Rightarrow\mu
   \andalso a\ne b}
  {\Gamma\vdash p.a\Rightarrow\mu}
\end{minipage}

\plateheader{Capture-set and member-signature well-formedness}
 {\Gamma\vdash C\ \wf\qquad\Gamma\vdash\mu\ \wf}
\noindent
\begin{minipage}[t]{0.48\linewidth}
\noindent\textit{Capture sets}\par\nobreak
\begin{minipage}[t]{0.46\linewidth}
\infax[\ruledef{wfc-empty}]
  {\Gamma\vdash\cempty\ \wf}

\infrule[\ruledef{wfc-union}]
  {\Gamma\vdash C\ \wf\quad\Gamma\vdash D\ \wf}
  {\Gamma\vdash C\cup D\ \wf}
\end{minipage}\hfill
\begin{minipage}[t]{0.51\linewidth}
\infrule[\ruledef{wfc-path}]
  {\Gamma\vdash p\Rightarrow T}
  {\Gamma\vdash\cset p\ \wf}

\infrule[\ruledef{wfc-select}]
  {\Gamma\vdash p.a\Rightarrow C..D
   \\
   \Gamma\vdash C<:D}
  {\Gamma\vdash p.a\ \wf}
\end{minipage}
\end{minipage}\hfill
\begin{minipage}[t]{0.50\linewidth}
\noindent\textit{Member signatures}\par\nobreak
\begin{minipage}[t]{0.31\linewidth}
\infrule[\ruledef{wfm-term}]
  {\Gamma\vdash T\ \wf}
  {\Gamma\vdash T\ \wf}
\end{minipage}\hfill
\begin{minipage}[t]{0.66\linewidth}
\infrule[\ruledef{wfm-type}]
  {\Gamma\vdash S\ \wf\quad\Gamma\vdash R\ \wf
   \quad\Gamma\vdash S<:R}
  {\Gamma\vdash S..R\ \wf}

\infrule[\ruledef{wfm-capture}]
  {\Gamma\vdash C\ \wf\quad\Gamma\vdash D\ \wf
   \quad\Gamma\vdash C<:D}
  {\Gamma\vdash C..D\ \wf}
\end{minipage}
\end{minipage}

\plateheader{Qualified-type and shape well-formedness}
  {\Gamma\vdash T\ \wf\quad\Gamma\vdash S\ \wf}
\noindent
\begin{minipage}[t]{0.31\linewidth}
\infrule[\ruledef{wft-capt}]
  {\Gamma\vdash C\ \wf\andalso\Gamma\vdash S\ \wf}
  {\Gamma\vdash\capt{S}{C}\ \wf}

\infax[\ruledef{wfs-bot}]
  {\Gamma\vdash\bot\ \wf}

\infax[\ruledef{wfs-top}]
  {\Gamma\vdash\top\ \wf}
\end{minipage}\hfill
\begin{minipage}[t]{0.32\linewidth}
\infrule[\ruledef{wfs-single}]
  {\Gamma\vdash p\Rightarrow T}
  {\Gamma\vdash\sing p\ \wf}

\infrule[\ruledef{wfs-select}]
  {\Gamma\vdash p.a\Rightarrow S..R
   \andalso\Gamma\vdash S<:R}
  {\Gamma\vdash p.a\ \wf}
\end{minipage}\hfill
\begin{minipage}[t]{0.34\linewidth}
\infrule[\ruledef{wfs-fun}]
  {\Gamma\vdash T\ \wf
   \andalso\Gamma,x:T\vdash U\ \wf}
  {\Gamma\vdash(x:T)\to U\ \wf}

\infrule[\ruledef{wfs-pair}]
  {\Gamma\vdash T\ \wf
   \andalso\Gamma,x:T\vdash\mu\ \wf}
  {\Gamma\vdash\pairshape{x}{T}{a}{\mu}\ \wf}
\end{minipage}

\caption{Source syntax and static semantics: source syntax, path synthesis,
and well-formedness judgments.}
\Description{The complete source grammar, path-synthesis judgment, and all
well-formedness judgments.}
\label{fig:syntax-wf}
\label{fig:path-synthesis}
\label{fig:well-formedness}
\end{figure}

\clearpage
\begin{figure}[H]
\ContinuedFloat
\plateheader{Subcapturing}{\Gamma\vdash C<:D}
\noindent
\begin{minipage}[t]{0.31\linewidth}
\infax[\ruledef{sc-refl}]
  {\Gamma\vdash C<:C}

\infrule[\ruledef{sc-trans}]
  {\Gamma\vdash C<:D
   \andalso
   \Gamma\vdash D<:C'}
  {\Gamma\vdash C<:C'}

\infax[\ruledef{sc-empty}]
  {\Gamma\vdash\cempty<:C}

\infax[\ruledef{sc-union-left}]
  {\Gamma\vdash C<:C\cup D}
\end{minipage}\hfill
\begin{minipage}[t]{0.32\linewidth}
\infax[\ruledef{sc-union-right}]
  {\Gamma\vdash D<:C\cup D}

\infrule[\ruledef{sc-union-elim}]
  {\Gamma\vdash C<:C'
   \andalso
   \Gamma\vdash D<:C'}
  {\Gamma\vdash C\cup D<:C'}

\infrule[\ruledef{sc-path}]
  {\Gamma\vdash p\Rightarrow\capt{S}{C}}
  {\Gamma\vdash\cset p<:C}

\infrule[\ruledef{sc-alias}]
  {\Gamma\vdash p\Rightarrow\capt{\sing q}{C}}
  {\Gamma\vdash\cset q<:\cset p}
\end{minipage}\hfill
\begin{minipage}[t]{0.34\linewidth}
\infrule[\ruledef{sc-fst-root}]
  {\Gamma\vdash p.1\Rightarrow T}
  {\Gamma\vdash\cset{p.1}<:\cset p}

\infrule[\ruledef{sc-sel-root}]
  {\Gamma\vdash p.a\Rightarrow T}
  {\Gamma\vdash\cset{p.a}<:\cset p}

\infrule[\ruledef{sc-select-lower}]
  {\Gamma\vdash p.a\Rightarrow L..U
   \andalso
   \Gamma\vdash L<:U}
  {\Gamma\vdash L<:p.a}

\infrule[\ruledef{sc-select-upper}]
  {\Gamma\vdash p.a\Rightarrow L..U
   \andalso
   \Gamma\vdash L<:U}
  {\Gamma\vdash p.a<:U}
\end{minipage}

\plateheader{Qualified-type subtyping}{\Gamma\vdash T<:U}
\noindent
\begin{minipage}[t]{0.30\linewidth}
\infax[\ruledef{st-refl}]
  {\Gamma\vdash T<:T}
\end{minipage}\hfill
\begin{minipage}[t]{0.31\linewidth}
\infrule[\ruledef{st-trans}]
  {\Gamma\vdash T<:U
   \andalso
   \Gamma\vdash U<:T'}
  {\Gamma\vdash T<:T'}
\end{minipage}\hfill
\begin{minipage}[t]{0.35\linewidth}
\infrule[\ruledef{st-capt}]
  {\Gamma\vdash C<:D
   \andalso
   \Gamma\vdash S<:R}
  {\Gamma\vdash\capt{S}{C}<:\capt{R}{D}}
\end{minipage}

\plateheader{Shape subtyping}{\Gamma\vdash S<:R}
\noindent
\begin{minipage}[t]{0.31\linewidth}
\infax[\ruledef{ss-refl}]
  {\Gamma\vdash S<:S}

\infrule[\ruledef{ss-trans}]
  {\Gamma\vdash S<:R
   \andalso
   \Gamma\vdash R<:S'}
  {\Gamma\vdash S<:S'}

\infax[\ruledef{ss-bot}]
  {\Gamma\vdash\bot<:S}

\infax[\ruledef{ss-top}]
  {\Gamma\vdash S<:\top}
\end{minipage}\hfill
\begin{minipage}[t]{0.32\linewidth}
\infrule[\ruledef{ss-widen}]
  {\Gamma\vdash p\Rightarrow\capt{S}{C}}
  {\Gamma\vdash\sing p<:S}
\infrule[\ruledef{ss-alias}]
  {\Gamma\vdash p\Rightarrow\capt{\sing q}{C}}
  {\Gamma\vdash\sing q<:\sing p}

\infrule[\ruledef{ss-select-lower}]
  {\Gamma\vdash p.a\Rightarrow S..R
   \andalso
   \Gamma\vdash S<:R}
  {\Gamma\vdash S<:p.a}
\end{minipage}\hfill
\begin{minipage}[t]{0.34\linewidth}
\infrule[\ruledef{ss-select-upper}]
  {\Gamma\vdash p.a\Rightarrow S..R
   \andalso
   \Gamma\vdash S<:R}
  {\Gamma\vdash p.a<:R}

\infrule[\ruledef{ss-fun}]
  {\Gamma\vdash T'<:T
   \andalso
   \Gamma,x:T'\vdash U<:U'}
  {\Gamma\vdash(x:T)\to U<:(x:T')\to U'}

\infrule[\ruledef{ss-pair}]
  {\Gamma\vdash T<:T'
   \andalso
   \Gamma,x:T\vdash\mu<:\mu'}
  {\Gamma\vdash
    \pairshape{x}{T}{a}{\mu}<:
    \pairshape{x}{T'}{a}{\mu'}}
\end{minipage}

\plateheader{Member-signature subtyping}{\Gamma\vdash\mu<:\nu}
\noindent
\begin{minipage}[t]{0.47\linewidth}
\infax[\ruledef{sm-refl}]
  {\Gamma\vdash\mu<:\mu}

\infrule[\ruledef{sm-trans}]
  {\Gamma\vdash\mu<:\nu
   \andalso
   \Gamma\vdash\nu<:\mu'}
  {\Gamma\vdash\mu<:\mu'}

\infrule[\ruledef{sm-term}]
  {\Gamma\vdash T<:U}
  {\Gamma\vdash T<:U}
\end{minipage}\hfill
\begin{minipage}[t]{0.51\linewidth}
\infrule[\ruledef{sm-type}]
  {\Gamma\vdash S'<:S
   \andalso
   \Gamma\vdash R<:R'
   \andalso
   \Gamma\vdash S<:R}
  {\Gamma\vdash S..R<:S'..R'}

\infrule[\ruledef{sm-capture}]
  {\Gamma\vdash C'<:C
   \andalso
   \Gamma\vdash D<:D'
   \andalso
   \Gamma\vdash C<:D}
  {\Gamma\vdash C..D<:C'..D'}
\end{minipage}

\caption{Static semantics (continued): subcapturing and all subtyping
judgments.}
\Description{The complete subcapturing, qualified-type, shape, and
member-signature subtyping judgments.}
\label{fig:subcapturing}
\label{fig:subtyping}
\label{fig:member-subtyping}
\end{figure}

\clearpage
\begin{figure}[H]
\ContinuedFloat

\plateheader{Term typing}{\Gamma\vdash t:T\uses C}
\noindent
\begin{minipage}[t]{0.47\linewidth}
\infrule[\ruledef{t-path}]
  {\Gamma\vdash p\Rightarrow T}
  {\Gamma\vdash p:\capt{\sing p}{\cset p}\uses\cset p}

\infrule[\ruledef{t-abs}]
  {\Gamma,x:T\vdash t:U\uses C^{\uparrow}\cup\cset x
   \\
   \Gamma\vdash T\ \wf
   \andalso
   \Gamma\vdash C\ \wf}
  {\Gamma\vdash\lambda(x:T).t:
    \capt{(x:T)\to U}{C}\uses\cempty}

\infrule[\ruledef{t-app}]
  {\Gamma\vdash p:\capt{(x:T)\to U}{C}\uses C_p
   \\
   \Gamma\vdash q:T\uses C_q}
  {\Gamma\vdash p\,q:\subst{U}{q}{x}\uses C_p\cup C_q}

\infax[\ruledef{t-pair}]
  {\begin{gathered}
   \Gamma\vdash\pairval{y}{a}{z}:\\[-2pt]
   \capt{\pairshape{x}{\capt{\sing y}{\cset y}}{a}
     {\capt{\sing{z^{\uparrow}}}{\cset{z^{\uparrow}}}}}
     {\cset y\cup\cset z}\uses\cempty
   \end{gathered}}
\end{minipage}\hfill
\begin{minipage}[t]{0.51\linewidth}
\infrule[\ruledef{t-type-pair}]
  {\Gamma\vdash S\ \wf}
  {\Gamma\vdash\pairval{y}{a}{S}:
   \capt{\pairshape{x}{\capt{\sing y}{\cset y}}{a}
     {S^{\uparrow}..S^{\uparrow}}}{\cset y}\uses\cempty}

\infrule[\ruledef{t-capture-pair}]
  {\Gamma\vdash C\ \wf}
  {\Gamma\vdash\pairval{y}{a}{C}:
   \capt{\pairshape{x}{\capt{\sing y}{\cset y}}{a}
     {C^{\uparrow}..C^{\uparrow}}}{\cset y}\uses\cempty}

\infrule[\ruledef{t-let}]
  {\Gamma\vdash t:T\uses C
   \andalso
   \Gamma,x:T\vdash u:U^{\uparrow}\uses C^{\uparrow}
   \\
   \Gamma\vdash U\ \wf
   \andalso
   \Gamma\vdash C\ \wf}
  {\Gamma\vdash\mathbf{let}\ x=t\ \mathbf{in}\ u:U\uses C}

\infrule[\ruledef{t-sub}]
  {\Gamma\vdash t:T\uses C
   \andalso
   \Gamma\vdash T<:U
   \andalso
   \Gamma\vdash C<:D
   \\
   \Gamma\vdash U\ \wf
   \andalso
   \Gamma\vdash D\ \wf}
  {\Gamma\vdash t:U\uses D}
\end{minipage}

\noindent
\begin{minipage}[t]{0.59\linewidth}
\plateheader{Runtime syntax}{}
\[
\begin{array}{rcl@{\qquad}l}
r &::=& \loc x\mid\typeref S\mid\capref C & \text{referent}\cr
\sigma &::=& \cempty\mid\sigma\mathbin{\cdot}(x\mapsto v)
  & \text{store}\cr
k &::=& [\,]\mid
  (\mathbf{let}\ x=[\,]\ \mathbf{in}\ u)::k
  & \text{continuation}\cr
s &::=& \langle\sigma,k,t\rangle & \text{state}
\end{array}
\]
\[
\refdef z=\loc z\qquad
\refdef S=\typeref S\qquad
\refdef C=\capref C.
\]
\end{minipage}\hfill
\begin{minipage}[t]{0.38\linewidth}
\plateheader{Final state}{\mathsf{final}(s)}
\infax[\ruledef{f-location}]
  {\mathsf{final}(\langle\sigma,[\,],x\rangle)}

\infrule[\ruledef{f-value}]
  {v\ \mathsf{value}}
  {\mathsf{final}(\langle\sigma,[\,],v\rangle)}
\end{minipage}

\plateheader{Generalized path resolution}{\sigma\vdash p\resolve r}
\noindent
\begin{minipage}[t]{0.47\linewidth}
\infax[\ruledef{r-var}]
  {\sigma\vdash x\resolve\loc x}

\infrule[\ruledef{r-fst}]
  {\sigma\vdash p\resolve\loc z
   \andalso
   \sigma(z)=\pairval{y}{a}{\delta}}
  {\sigma\vdash p.1\resolve\loc y}
\end{minipage}\hfill
\begin{minipage}[t]{0.51\linewidth}
\infrule[\ruledef{r-sel}]
  {\sigma\vdash p\resolve\loc z
   \andalso
   \sigma(z)=\pairval{y}{a}{\delta}}
  {\sigma\vdash p.a\resolve\refdef\delta}

\infrule[\ruledef{r-miss}]
  {\sigma\vdash p\resolve\loc z
   \andalso
   \sigma(z)=\pairval{y}{b}{\delta}
   \\
   a\ne b
   \andalso
   \sigma\vdash y.a\resolve r}
  {\sigma\vdash p.a\resolve r}
\end{minipage}

\plateheader{Machine transition}
 {\langle\sigma,k,t\rangle\sstep\langle\sigma',k',t'\rangle}
\noindent
\begin{minipage}[t]{0.47\linewidth}
\infrule[\ruledef{e-app}]
  {\sigma\vdash p\resolve\loc f
   \andalso
   \sigma\vdash q\resolve\loc y
   \\
   \sigma(f)=\lambda(x:T).t}
  {\langle\sigma,k,p\,q\rangle\sstep
   \langle\sigma,k,\subst{t}{y}{x}\rangle}

\infrule[\ruledef{e-path}]
  {\sigma\vdash p\resolve\loc x
   \andalso p\ \mathsf{not\ a\ variable}}
  {\langle\sigma,k,p\rangle\sstep
   \langle\sigma,k,x\rangle}

\infax[\ruledef{e-let}]
  {\langle\sigma,k,\mathbf{let}\ x=t\ \mathbf{in}\ u\rangle
   \sstep
   \langle\sigma,(\mathbf{let}\ x=[\,]\ \mathbf{in}\ u)::k,t\rangle}
\end{minipage}\hfill
\begin{minipage}[t]{0.51\linewidth}
\infax[\ruledef{e-return}]
  {\langle\sigma,
      (\mathbf{let}\ x=[\,]\ \mathbf{in}\ u)::k,y\rangle
   \sstep
   \langle\sigma,k,\subst{u}{y}{x}\rangle}

\infrule[\ruledef{e-allocate}]
  {v\ \mathsf{value}
   \andalso z\notin\mathsf{dom}(\sigma)}
  {\langle\sigma,
      (\mathbf{let}\ x=[\,]\ \mathbf{in}\ u)::k,v\rangle
   \sstep\\
   \langle\sigma\mathbin{\cdot}(z\mapsto v),k,
     \subst{u}{z}{x}\rangle}
\end{minipage}

\caption{Static and dynamic semantics (continued): term typing, runtime
syntax, path resolution, and machine transitions.}
\Description{The complete term-typing judgment, runtime grammar,
final-state judgment, generalized path-resolution judgment, and CK-machine
transition relation.}
\label{fig:term-typing}
\label{fig:typing-runtime}
\label{fig:runtime}
\end{figure}

\endgroup
\clearpage

Path synthesis \(\Gamma\vdash p\Rightarrow\mu\) is precise and has no
subsumption. The outer qualifier of a pair is irrelevant to lookup. Matching
selection opens the dependent member with the receiver's first projection.
When labels differ, selection continues through the first component; this
represents a right-nested record without a separate record calculus.

There are four subtyping judgments. Subcapturing
\(\Gamma\vdash C<:D\) compares capture sets. Qualified-type and shape
subtyping are written \(\Gamma\vdash T<:U\) and
\(\Gamma\vdash S<:R\); the syntactic categories disambiguate them.
Generalized signature subtyping uses the same notation. Interval
well-formedness and the type- and capture-selection subtyping rules carry
consistency premises. A qualified type may also be regarded as a generalized
term signature. Since the grammar presents this inclusion without a
constructor, the premise and conclusion of the corresponding
well-formedness and subtyping rules have the same typography but belong to
different judgment families.

The abstraction rule deserves one point of notation. If \(C\) is formed in
\(\Gamma\), then \(C^{\uparrow}\) is its weakening to
\(\Gamma,x:T\). The premise use set
\(C^{\uparrow}\cup\cset x\) contains no free occurrence of the new binder
after abstraction. Subsumption and the proper-path root rules allow any path
rooted at \(x\), such as \(x.1.a\), to be widened to \(\cset x\). Let
typing similarly requires the body's result type and use bound to be
formed in the outer context; subsumption must eliminate the local root before
the rule applies.

We write
\(\mathsf{initial}(t)=\langle\varnothing,[\,],t\rangle\).  A store is an
immutable sequence of allocated values; \(\mathsf{dom}(\sigma)\) is exactly
its set of allocated locations.  In particular, every runtime variable names
a store cell.  Allocation chooses a fresh location and transports the old
store, continuation, and term into the extended location scope.  The named
presentation of \ruleref{e-allocate} leaves this transport implicit.

Call a generalized signature bound-consistent in \(\Gamma\) when it is a
term signature, or when it is an interval whose lower bound is below its
upper bound in the corresponding subtype judgment.

\begin{lemma}[Bound consistency]
\label{lem:bound-consistency}
Every well-formed generalized signature is bound-consistent.  If
\(\Gamma\vdash\mu<:\nu\) and \(\mu\) is bound-consistent, then \(\nu\) is
bound-consistent.
\end{lemma}

The well-formedness result follows immediately from the interval rules.  For
signature subtyping, the type-interval case chains
\(S'<:S<:R<:R'\), and the capture-interval case chains
\(C'<:C<:D<:D'\); reflexivity and transitivity preserve the property.
The mechanization does not yet include a full context-formation and typing
regularity theorem.  That stronger result requires the opening and path
substitution metatheory needed to recover well-formed synthesized signatures
from arbitrary typing derivations.

\begin{lemma}[Resolution determinism]
\label{lem:resolution-determinism}
If \(\sigma\vdash p\resolve r\) and
\(\sigma\vdash p\resolve r'\), then \(r=r'\).
\end{lemma}

The proof follows the first resolution derivation.  In the two selection
cases, immutability gives a unique store binding, and the induction
hypothesis determines the tail referent.  Generalized resolution is used in
the static and semantic arguments; the machine steps only when its result is
a location.

\subsection{Capture members and root contraction}

A capture definition packages a capture set behind a stable path.  For
example, after allocating
\[
  r=\pairval{y}{a}{W},
\]
where \(W\) is a capture set, path synthesis exposes the exact interval
\(r.a\Rightarrow W..W\).  Subcapturing may subsequently forget precision by
placing \(r.a\) between wider declared bounds.  Capture selection therefore
provides capture-set abstraction without adding a separate binder form or
explicit capture polymorphism to terms.

The root rules apply only to singleton capture sets of proper term paths.
They express that using \(p.1\) or \(p.a\) requires access to the root
capability \(p\).  An abstract capture selection has no such ownership
interpretation.  This restriction is visible in the type of a concrete
capture pair: its outer qualifier is \(\cset y\), independent of the stored
definition \(W\).  If one added \(p.a<:\cset p\), the exact lower bound
\(W<:p.a\), together with \(\cset p<:\cset y\), would derive
\(W<:\cset y\) for an arbitrary stored capture set.  The qualifier would
then cease to bound the capabilities retained by the definition.

The two syntactic occurrences of paths in capture sets also remain distinct.
The singleton \(\cset p\) denotes the capability named by a term path; the
selection \(p.a\) denotes the capture set stored in a capture member.  This
distinction is used by both path synthesis and the capture-coercion rules.

\subsection{The general pair rule}

Rule \ruleref{ss-pair} compares dependent members under the source type of
the first component.  It supports ordinary covariance of the first component
without requiring the member signature to be independent of the binder.  Two
representative instances are
\[
\begin{array}{c}
 x:\top\vdash \bot<:\sing x
 \qquad
 x:\top\vdash \sing x<:\top
 \\[-1pt]
 \hline
 \pairshape{x}{\top}{a}{\sing x..\sing x}
 <: \pairshape{x}{\top}{a}{\bot..\top}
\end{array}
\qquad
\begin{array}{c}
 x:\top\vdash \cempty<:\cset x
 \qquad
 x:\top\vdash \cset x<:\cset x
 \\[-1pt]
 \hline
 \pairshape{x}{\top}{a}{\cset x..\cset x}
 <: \pairshape{x}{\top}{a}{\cempty..\cset x}.
\end{array}
\]
The mechanization exercises both derivations in closed executions that
allocate the coerced pair, thereby forcing the semantic pair action rather
than merely constructing a source derivation.  A separate capture-selection
execution forces both the lower- and upper-bound selection coercions.  The
examples are useful because the member premises genuinely mention \(x\), even
though the first-component types on both sides are \(\top\).

\section{Qualifier-aware store interpretation}
\label{sec:semantics}

The preservation argument translates source subtyping derivations into
store-indexed coercions.  Coercion application maps semantic evidence for a
source type to evidence for its target type, so preservation never needs to
invert a derivation ending in transitivity.  The interpretation also retains
the two capture-related indices of source typing: the introduction qualifier
of every stored value and the use set of every running term.

\subsection{Runtime path equality and conversion}

The store equates paths that resolve to the same referent.  We write
\(p\equiv_\sigma q\) for the least equivalence and path congruence generated
by co-resolution.  In \ruleref{peq-cong}, \(p_0\) is formed under one
additional variable \(x\).

\begin{figure}[H]
\begingroup
\rulesetup
\plateheader{Runtime path equality}{p\equiv_\sigma q}
\noindent
\begin{minipage}[t]{0.47\linewidth}
\infax[\ruledef{peq-refl}]
  {p\equiv_\sigma p}

\infrule[\ruledef{peq-symm}]
  {p\equiv_\sigma q}
  {q\equiv_\sigma p}

\infrule[\ruledef{peq-trans}]
  {p\equiv_\sigma q\andalso q\equiv_\sigma p'}
  {p\equiv_\sigma p'}
\end{minipage}\hfill
\begin{minipage}[t]{0.51\linewidth}
\infrule[\ruledef{peq-coresolve}]
  {\sigma\vdash p\resolve r
   \andalso
   \sigma\vdash q\resolve r}
  {p\equiv_\sigma q}

\infrule[\ruledef{peq-cong}]
  {p\equiv_\sigma q}
  {\subst{p_0}{p}{x}\equiv_\sigma\subst{p_0}{q}{x}}
\end{minipage}
\endgroup
\caption{Store-induced equality of paths.}
\Description{All rules defining runtime path equality.}
\label{fig:runtime-equality}
\end{figure}

Runtime conversion \(X\simeq_\sigma Y\) is defined for capture sets,
qualified types, shapes, and member signatures.  For each category it is the
least reflexive, symmetric, transitive, constructor-compatible relation
containing
\[
  \frac{p\equiv_\sigma q}
       {\subst{X}{p}{x}\simeq_\sigma\subst{X}{q}{x}}.
\]
Conversion preserves the outer constructor and member kind.  Beneath a
binder, ambient equalities are weakened and the bound variable is related
only to itself; we write the resulting conversion as
\(X\simeq^x_\sigma Y\).  Opening it with runtime-equal paths yields ordinary
conversion.

\subsection{Annotated stores and semantic judgments}

An annotated store \(\Omega\) has the same spine as \(\sigma\) and records an
introduction qualifier beside each stored value:
\[
  \Omega ::= [\,] \mid \Omega\mathbin{\cdot}(v@Q).
\]
We write \(\Omega(x)=v@Q\) for lookup.  All locations in \(v\) and \(Q\)
refer to earlier cells, as they do in the underlying store.  The judgment
\(\Omega\valid\sigma\), defined below together with value evidence, ensures
that each annotation is justified by the value rule that produced the same
stored value.

Let \(\qual(\capt S C)=C\).  The core semantic judgments are
\[
\begin{array}{rcll}
 \Omega;\sigma\vdash C\semleq D
   &&& \text{semantic subcapturing},\\
 \Omega;\sigma\vdash T\tcoerce U
   &&& \text{qualified-type coercion},\\
 \Omega;\sigma\vdash S\scoerce R
   &&& \text{shape coercion},\\
 \Omega;\sigma\vdash\mu\coerce\nu
   &&& \text{member coercion},\\
 \Omega;\sigma\models x:T@Q
   &&& \text{location inhabitation with introduction qualifier }Q,\\
 \Omega;\sigma\models r:\mu
   &&& \text{generalized-referent realization}.
\end{array}
\]
They are mutually defined.  Source premises are closed by a valuation
\(\rho:\mathsf{dom}(\Gamma)\to\mathsf{dom}(\sigma)\).  The notation
\(\rho(X)\) replaces the free variables of \(X\) with their locations while
leaving a displayed binder in place.  Valuations need not be injective.

\clearpage
\begingroup
\setlength{\abovedisplayskip}{5pt plus 1pt minus 1pt}
\setlength{\belowdisplayskip}{5pt plus 1pt minus 1pt}
\setlength{\intextsep}{5pt}
\captionsetup{skip=6pt}
\rulesetup

\begin{figure}[H]
\plateheader{Semantic subcapturing}{\Omega;\sigma\vdash C\semleq D}
\noindent
\begin{minipage}[t]{0.31\linewidth}
\infrule[\ruledef{cap-source}]
  {\rho\models^\Omega_\sigma\Gamma
   \\
   \Gamma\vdash C<:D}
  {\Omega;\sigma\vdash\rho(C)\semleq\rho(D)}

\infax[\ruledef{cap-refl}]
  {\Omega;\sigma\vdash C\semleq C}

\infrule[\ruledef{cap-trans}]
  {\Omega;\sigma\vdash C\semleq D
   \\
   \Omega;\sigma\vdash D\semleq E}
  {\Omega;\sigma\vdash C\semleq E}

\infrule[\ruledef{cap-conv}]
  {C\simeq_\sigma D}
  {\Omega;\sigma\vdash C\semleq D}

\infax[\ruledef{cap-empty}]
  {\Omega;\sigma\vdash\cempty\semleq C}
\end{minipage}\hfill
\begin{minipage}[t]{0.31\linewidth}
\infax[\ruledef{cap-union-left}]
  {\Omega;\sigma\vdash C\semleq C\cup D}

\infax[\ruledef{cap-union-right}]
  {\Omega;\sigma\vdash D\semleq C\cup D}

\infrule[\ruledef{cap-union-elim}]
  {\Omega;\sigma\vdash C\semleq E
   \\
   \Omega;\sigma\vdash D\semleq E}
  {\Omega;\sigma\vdash C\cup D\semleq E}

\infrule[\ruledef{cap-alias}]
  {\sigma\vdash p\resolve\loc x
   \\
   \sigma\vdash q\resolve\loc x}
  {\Omega;\sigma\vdash\cset q\semleq\cset p}

\infrule[\ruledef{cap-fold}]
  {\sigma\vdash p\resolve\loc x
   \\
   \Omega(x)=v@Q}
  {\Omega;\sigma\vdash Q\semleq\cset p}
\end{minipage}\hfill
\begin{minipage}[t]{0.34\linewidth}
\infrule[\ruledef{cap-fst-root}]
  {\sigma\vdash p.1\resolve\loc x}
  {\Omega;\sigma\vdash\cset{p.1}\semleq\cset p}

\infrule[\ruledef{cap-sel-root}]
  {\sigma\vdash p.a\resolve\loc x}
  {\Omega;\sigma\vdash\cset{p.a}\semleq\cset p}

\infrule[\ruledef{cap-select-lower}]
  {\sigma\vdash p.a\resolve\capref W
   \\
   \Omega;\sigma\vdash L\semleq W}
  {\Omega;\sigma\vdash L\semleq p.a}

\infrule[\ruledef{cap-select-upper}]
  {\sigma\vdash p.a\resolve\capref W
   \\
   \Omega;\sigma\vdash W\semleq U}
  {\Omega;\sigma\vdash p.a\semleq U}
\end{minipage}

\caption{The complete semantic subcapturing relation.}
\Description{All rules of world-indexed semantic subcapturing.}
\label{fig:semantic-subcapturing}
\end{figure}

\begin{figure}[H]
\ContinuedFloat
\plateheader{Qualified-type coercion}{\Omega;\sigma\vdash T\tcoerce U}
\noindent
\begin{minipage}[t]{0.22\linewidth}
\infax[\ruledef{tc-refl}]
  {\Omega;\sigma\vdash T\tcoerce T}
\end{minipage}\hfill
\begin{minipage}[t]{0.27\linewidth}
\infrule[\ruledef{tc-trans}]
  {\Omega;\sigma\vdash T\tcoerce U
   \\
   \Omega;\sigma\vdash U\tcoerce T'}
  {\Omega;\sigma\vdash T\tcoerce T'}
\end{minipage}\hfill
\begin{minipage}[t]{0.20\linewidth}
\infrule[\ruledef{tc-conv}]
  {T\simeq_\sigma U}
  {\Omega;\sigma\vdash T\tcoerce U}
\end{minipage}\hfill
\begin{minipage}[t]{0.27\linewidth}
\infrule[\ruledef{tc-capt}]
  {\Omega;\sigma\vdash C\semleq D
   \\
   \Omega;\sigma\vdash S\scoerce R}
  {\Omega;\sigma\vdash\capt S C\tcoerce\capt R D}
\end{minipage}

\caption{The complete qualified-type coercion.}
\Description{All rules for qualifier-aware type coercion.}
\label{fig:type-coercion}
\end{figure}

\begin{figure}[H]
\ContinuedFloat
\plateheader{Shape coercion}{\Omega;\sigma\vdash S\scoerce R}
\noindent
\begin{minipage}[t]{0.30\linewidth}
\infax[\ruledef{shc-refl}]
  {\Omega;\sigma\vdash S\scoerce S}

\infrule[\ruledef{shc-trans}]
  {\Omega;\sigma\vdash S\scoerce R
   \\
   \Omega;\sigma\vdash R\scoerce S'}
  {\Omega;\sigma\vdash S\scoerce S'}

\infrule[\ruledef{shc-conv}]
  {S\simeq_\sigma R}
  {\Omega;\sigma\vdash S\scoerce R}

\infax[\ruledef{shc-bot}]
  {\Omega;\sigma\vdash\bot\scoerce S}

\infax[\ruledef{shc-top}]
  {\Omega;\sigma\vdash S\scoerce\top}
\end{minipage}\hfill
\begin{minipage}[t]{0.30\linewidth}
\infrule[\ruledef{shc-widen}]
  {\sigma\vdash p\resolve\loc x
   \\
   \Omega;\sigma\models x:\capt S C@Q}
  {\Omega;\sigma\vdash\sing p\scoerce S}

\infrule[\ruledef{shc-alias}]
  {\sigma\vdash p\resolve\loc x
   \\
   \sigma\vdash q\resolve\loc x}
  {\Omega;\sigma\vdash\sing q\scoerce\sing p}

\infrule[\ruledef{shc-select-lower}]
  {\sigma\vdash p.a\resolve\typeref W
   \\
   \Omega;\sigma\vdash L\scoerce W}
  {\Omega;\sigma\vdash L\scoerce p.a}

\infrule[\ruledef{shc-select-upper}]
  {\sigma\vdash p.a\resolve\typeref W
   \\
   \Omega;\sigma\vdash W\scoerce U}
  {\Omega;\sigma\vdash p.a\scoerce U}
\end{minipage}\hfill
\begin{minipage}[t]{0.36\linewidth}
\infrule[\ruledef{shc-fun}]
  {\Omega;\sigma\vdash T'\tcoerce T
   \\
   \Omega;\sigma;x:T'\vdash
      U\rightsquigarrow_{\mathsf d}U'}
  {\Omega;\sigma\vdash
    (x:T)\to U\scoerce(x:T')\to U'}

\infrule[\ruledef{shc-pair}]
  {\Omega;\sigma\vdash T\tcoerce T'
   \\
   \mathsf{Member}_{\Omega;\sigma}
      (x:T\vdash\mu<:\mu')}
  {\Omega;\sigma\vdash
    \pairshape{x}{T}{a}{\mu}\scoerce
    \pairshape{x}{T'}{a}{\mu'}}
\end{minipage}

\caption{The complete shape coercion.}
\Description{All rules for qualifier-aware shape coercion.}
\label{fig:shape-coercion}
\end{figure}

\begin{figure}[H]
\ContinuedFloat
\plateheader{Generalized member coercion}
  {\Omega;\sigma\vdash\mu\coerce\nu}
\noindent
\begin{minipage}[t]{0.31\linewidth}
\infax[\ruledef{mc-refl}]
  {\Omega;\sigma\vdash\mu\coerce\mu}

\infrule[\ruledef{mc-trans}]
  {\Omega;\sigma\vdash\mu\coerce\nu
   \\
   \Omega;\sigma\vdash\nu\coerce\mu'}
  {\Omega;\sigma\vdash\mu\coerce\mu'}
\end{minipage}\hfill
\begin{minipage}[t]{0.31\linewidth}
\infrule[\ruledef{mc-conv}]
  {\mu\simeq_\sigma\nu}
  {\Omega;\sigma\vdash\mu\coerce\nu}

\infrule[\ruledef{mc-term}]
  {\Omega;\sigma\vdash T\tcoerce U}
  {\Omega;\sigma\vdash T\coerce U}
\end{minipage}\hfill
\begin{minipage}[t]{0.34\linewidth}
\infrule[\ruledef{mc-type}]
  {\Omega;\sigma\vdash L'\scoerce L
   \andalso
   \Omega;\sigma\vdash U\scoerce U'}
  {\Omega;\sigma\vdash L..U\coerce L'..U'}

\infrule[\ruledef{mc-capture}]
  {\Omega;\sigma\vdash C'\semleq C
   \andalso
   \Omega;\sigma\vdash D\semleq D'}
  {\Omega;\sigma\vdash C..D\coerce C'..D'}
\end{minipage}

\caption{The complete generalized member coercion.}
\Description{All rules for generalized member coercion.}
\label{fig:member-coercions}
\end{figure}

\begin{figure}[H]
\ContinuedFloat
\plateheader{Suspended body}
 {\mathsf{Body}_{\Omega;\sigma}(x:T\vdash t:U\uses C)}
\infrule[\ruledef{body-source}]
  {\rho\models^\Omega_\sigma\Gamma
   \andalso
   \Gamma,x:S\vdash t:U\uses C}
  {\mathsf{Body}_{\Omega;\sigma}
    (x:\rho(S)\vdash\rho(t):\rho(U)\uses\rho(C))}

\plateheader{Suspended member premise}
 {\mathsf{Member}_{\Omega;\sigma}(x:T\vdash\mu<:\nu)}
\infrule[\ruledef{member-source}]
  {\rho\models^\Omega_\sigma\Gamma
   \andalso
   \Gamma,x:S\vdash\mu<:\nu}
  {\mathsf{Member}_{\Omega;\sigma}
    (x:\rho(S)\vdash\rho(\mu)<:\rho(\nu))}

\plateheader{Deferred function-result coercion}
 {\Omega;\sigma;x:T\vdash U\rightsquigarrow_{\mathsf d}U'}
\noindent
\begin{minipage}[t]{0.47\linewidth}
\infax[\ruledef{dc-refl}]
  {\Omega;\sigma;x:T\vdash U\rightsquigarrow_{\mathsf d}U}

\infrule[\ruledef{dc-trans}]
  {\Omega;\sigma;x:T\vdash U\rightsquigarrow_{\mathsf d}U'
   \\
   \Omega;\sigma;x:T\vdash U'\rightsquigarrow_{\mathsf d}U''}
  {\Omega;\sigma;x:T\vdash U\rightsquigarrow_{\mathsf d}U''}

\infrule[\ruledef{dc-conv}]
  {U\simeq^x_\sigma U'}
  {\Omega;\sigma;x:T\vdash U\rightsquigarrow_{\mathsf d}U'}
\end{minipage}\hfill
\begin{minipage}[t]{0.51\linewidth}
\infrule[\ruledef{dc-narrow}]
  {\Omega;\sigma\vdash T'\tcoerce T
   \\
   \Omega;\sigma;x:T\vdash U\rightsquigarrow_{\mathsf d}U'}
  {\Omega;\sigma;x:T'\vdash U\rightsquigarrow_{\mathsf d}U'}

\infrule[\ruledef{dc-source}]
  {\rho\models^\Omega_\sigma\Gamma
   \andalso
   \Gamma,x:T\vdash U<:U'}
  {\Omega;\sigma;x:\rho(T)\vdash
    \rho(U)\rightsquigarrow_{\mathsf d}\rho(U')}
\end{minipage}

\caption{Suspended source premises and the complete deferred
function-result coercion.}
\Description{The source rules for suspended bodies and dependent members,
and all rules for coercions suspended beneath a function binder.}
\label{fig:deferred-coercions}
\end{figure}

\endgroup

Rule \ruleref{cap-source} embeds a source subcapturing derivation together
with the environment that justifies its path premises.  The remaining rules
are the closure operations needed by coercion action and reduction.  In
particular, \ruleref{cap-fold} states that a path naming a stored value
summarizes that value's introduction qualifier.  There is no inverse rule.
Abstract capture selections are related only through their stored interval
bounds.

Function coercions retain their codomain premise until an argument location
is available.  Pair coercions retain the member premise until the stored
first component is available.  Type and capture intervals are contravariant
in the lower bound and covariant in the upper bound.  In the three source
closure rules, \(\rho\) acts on variables from \(\Gamma\); the displayed
\(x\) remains bound.

\subsection{Qualifier-aware inhabitants and valid stores}

Environment satisfaction, location inhabitation, referent realization, and
value evidence are shown in the next two figures.  Location evidence carries
an explicit introduction qualifier \(Q\).  Every location rule retrieves the
same \(Q\) from \(\Omega\) and proves \(Q\semleq C\), where \(C\) is the
outer qualifier of the assigned type.

\begingroup
\setlength{\abovedisplayskip}{5pt plus 1pt minus 1pt}
\setlength{\belowdisplayskip}{5pt plus 1pt minus 1pt}
\setlength{\intextsep}{5pt}
\captionsetup{skip=6pt}
\rulesetup

\begin{figure}[H]
\plateheader{Environment satisfaction}{\rho\models^\Omega_\sigma\Gamma}
\infrule[\ruledef{sem-env}]
  {\forall x\in\mathsf{dom}(\Gamma).\;
   \exists Q.\;
   \Omega;\sigma\models\rho(x):\rho(\Gamma(x))@Q}
  {\rho\models^\Omega_\sigma\Gamma}

\plateheader{Referent realization}{\Omega;\sigma\models r:\mu}
\noindent
\begin{minipage}[t]{0.31\linewidth}
\infrule[\ruledef{sem-ref-loc}]
  {\Omega;\sigma\models x:T@Q}
  {\Omega;\sigma\models\loc x:T}
\end{minipage}\hfill
\begin{minipage}[t]{0.32\linewidth}
\infrule[\ruledef{sem-ref-type}]
  {\Omega;\sigma\vdash L\scoerce W
   \\
   \Omega;\sigma\vdash W\scoerce U}
  {\Omega;\sigma\models\typeref W:L..U}
\end{minipage}\hfill
\begin{minipage}[t]{0.34\linewidth}
\infrule[\ruledef{sem-ref-capture}]
  {\Omega;\sigma\vdash L\semleq W
   \\
   \Omega;\sigma\vdash W\semleq U}
  {\Omega;\sigma\models\capref W:L..U}
\end{minipage}

\plateheader{Location inhabitation}{\Omega;\sigma\models x:T@Q}
\noindent
\begin{minipage}[t]{0.47\linewidth}
\infrule[\ruledef{sem-loc-top}]
  {\Omega(x)=v@Q
   \andalso
   \Omega;\sigma\vdash Q\semleq C}
  {\Omega;\sigma\models x:\capt\top C@Q}

\infrule[\ruledef{sem-loc-fun}]
  {\Omega(x)=\lambda(z:A).t@Q
   \\
   \mathsf{Body}_{\Omega;\sigma}
     (z:A\vdash t:B\uses Q^\uparrow\cup\cset z)
   \\
   \Omega;\sigma\vdash S\tcoerce A
   \andalso
   \Omega;\sigma;z:S\vdash B\rightsquigarrow_{\mathsf d}U
   \\
   \Omega;\sigma\vdash Q\semleq C}
  {\Omega;\sigma\models
    x:\capt{(z:S)\to U}{C}@Q}
\end{minipage}\hfill
\begin{minipage}[t]{0.51\linewidth}
\infrule[\ruledef{sem-loc-pair}]
  {\Omega(x)=\pairval{y}{a}{\delta}@Q
   \\
   \Omega;\sigma\models y:T@Q_y
   \\
   \Omega;\sigma\models
      \refdef\delta:\subst{\mu}{y}{z}
   \\
   \Omega;\sigma\vdash Q\semleq C}
  {\Omega;\sigma\models
    x:\capt{\pairshape{z}{T}{a}{\mu}}{C}@Q}

\infrule[\ruledef{sem-loc-single}]
  {\Omega(x)=v@Q
   \andalso
   \sigma\vdash p\resolve\loc x
   \\
   \Omega;\sigma\vdash Q\semleq C}
  {\Omega;\sigma\models x:\capt{\sing p}{C}@Q}

\infrule[\ruledef{sem-loc-select}]
  {\Omega(x)=v@Q
   \andalso
   \sigma\vdash p.a\resolve\typeref W
   \\
   \Omega;\sigma\models x:\capt W E@Q
   \andalso
   \Omega;\sigma\vdash Q\semleq C}
  {\Omega;\sigma\models x:\capt{p.a}{C}@Q}
\end{minipage}

\caption{Environment satisfaction, referent realization, and
qualifier-aware location inhabitation.}
\Description{All rules for semantic environments, generalized referents,
and location inhabitants.}
\label{fig:semantic-evidence}
\end{figure}

\endgroup

The introduction types of pair values are abbreviated as follows:
\[
\begin{array}{rcl}
I_{\mathsf{term}}(y,a,z)
 &=& \capt{
   \pairshape{x}{\capt{\sing y}{\cset y}}{a}
    {\capt{\sing{z^{\uparrow}}}{\cset{z^{\uparrow}}}}}
   {\cset y\cup\cset z},
\\[2pt]
I_{\mathsf{type}}(y,a,W)
 &=& \capt{
   \pairshape{x}{\capt{\sing y}{\cset y}}{a}
    {W^{\uparrow}..W^{\uparrow}}}
   {\cset y},
\\[2pt]
I_{\mathsf{capture}}(y,a,W)
 &=& \capt{
   \pairshape{x}{\capt{\sing y}{\cset y}}{a}
    {W^{\uparrow}..W^{\uparrow}}}
   {\cset y}.
\end{array}
\]
Here \(W\) is a shape in \(I_{\mathsf{type}}\) and a capture set in
\(I_{\mathsf{capture}}\).  The superscript \({}^{\uparrow}\) transports an
object beneath the pair binder without changing its free paths.

\begingroup
\setlength{\abovedisplayskip}{5pt plus 1pt minus 1pt}
\setlength{\belowdisplayskip}{5pt plus 1pt minus 1pt}
\setlength{\intextsep}{5pt}
\captionsetup{skip=6pt}
\rulesetup

\begin{figure}[H]
\plateheader{Value evidence}
 {\Omega;\sigma\vdash_{\mathsf v}v:T@Q}
\noindent
\begin{minipage}[t]{0.47\linewidth}
\infrule[\ruledef{ve-abs}]
  {\mathsf{Body}_{\Omega;\sigma}
     (x:A\vdash t:B\uses Q^\uparrow\cup\cset x)
   \\
   \Omega;\sigma\vdash
      \capt{(x:A)\to B}{Q}\tcoerce T}
  {\Omega;\sigma\vdash_{\mathsf v}
     \lambda(x:A).t:T@Q}

\infrule[\ruledef{ve-term-pair}]
  {\Omega;\sigma\vdash I_{\mathsf{term}}(y,a,z)\tcoerce T}
  {\Omega;\sigma\vdash_{\mathsf v}
     \pairval{y}{a}{z}:T@(\cset y\cup\cset z)}
\end{minipage}\hfill
\begin{minipage}[t]{0.51\linewidth}
\infrule[\ruledef{ve-type-pair}]
  {\Omega;\sigma\vdash I_{\mathsf{type}}(y,a,W)\tcoerce T}
  {\Omega;\sigma\vdash_{\mathsf v}
     \pairval{y}{a}{W}:T@\cset y}

\infrule[\ruledef{ve-capture-pair}]
  {\Omega;\sigma\vdash I_{\mathsf{capture}}(y,a,W)\tcoerce T}
  {\Omega;\sigma\vdash_{\mathsf v}
     \pairval{y}{a}{W}:T@\cset y}
\end{minipage}

\plateheader{Valid annotated store}{\Omega\valid\sigma}
\noindent
\begin{minipage}[t]{0.26\linewidth}
\infax[\ruledef{world-empty}]
  {[\,]\valid[\,]}
\end{minipage}\hfill
\begin{minipage}[t]{0.68\linewidth}
\infrule[\ruledef{world-allocate}]
  {\Omega\valid\sigma
   \andalso
   \Omega;\sigma\vdash_{\mathsf v}v:T@Q
   \andalso v\ \mathsf{value}}
  {\Omega\mathbin{\cdot}(v@Q)\valid
   \sigma\mathbin{\cdot}(z\mapsto v)}
\end{minipage}

\caption{Value evidence and validity of annotated stores.}
\Description{All value-evidence and annotated-store validity rules.}
\label{fig:value-world-evidence}
\end{figure}

\endgroup

The qualifier after \(@\) is fixed by the introduction rule.  It need not be
minimal or principal: the body derivation in \ruleref{ve-abs}, for example,
may already contain subsumption.  The coercion suffix records any further
subsumption above the value constructor.  Projecting the outer capture
component of that suffix gives
\[
  \Omega;\sigma\vdash_{\mathsf v}v:T@Q
  \quad\Longrightarrow\quad
  \Omega;\sigma\vdash Q\semleq\qual(T).
\]
Validity ties each annotation to value evidence for the same store entry.
Because stores are append-only, \ruleref{world-allocate} depends only on the
previous annotated store.  Allocation preserves all earlier evidence by
weakening locations uniformly.

\subsection{Compilation and coercion action}

Typed resolution and coercion compilation are stated under a satisfied
source environment:
\[
\begin{array}{rcl}
\Gamma\vdash p\Rightarrow\mu
 &\Longrightarrow&
 \exists r.\;\sigma\vdash\rho(p)\resolve r
       \ \land\ \Omega;\sigma\models r:\rho(\mu),
\\[3pt]
\Gamma\vdash C<:D
 &\Longrightarrow&
 \Omega;\sigma\vdash\rho(C)\semleq\rho(D),
\\
\Gamma\vdash T<:U
 &\Longrightarrow&
 \Omega;\sigma\vdash\rho(T)\tcoerce\rho(U),
\\
\Gamma\vdash S<:R
 &\Longrightarrow&
 \Omega;\sigma\vdash\rho(S)\scoerce\rho(R),
\\
\Gamma\vdash\mu<:\nu
 &\Longrightarrow&
 \Omega;\sigma\vdash\rho(\mu)\coerce\rho(\nu).
\end{array}
\]
Each construction follows its source derivation.  Transitivity builds an
explicit composition.  Function subtyping produces \ruleref{shc-fun} with a
deferred result coercion.  Pair subtyping produces \ruleref{shc-pair} with a
member closure.  Neither case requires transitivity inversion.

\begin{lemma}[Qualifier-preserving coercion action]
\label{lem:coercion-action}
If \(\Omega;\sigma\models r:\mu\) and
\(\Omega;\sigma\vdash\mu\coerce\nu\), then
\(\Omega;\sigma\models r:\nu\).  In particular,
\[
 \Omega;\sigma\models x:T@Q
 \quad\land\quad
 \Omega;\sigma\vdash T\tcoerce U
 \quad\Longrightarrow\quad
 \Omega;\sigma\models x:U@Q.
\]
\end{lemma}

Identity returns its argument, composition applies its components in order,
and runtime conversion transports evidence along store-induced path
equality.  Interval action composes the coercions on their bounds.  For
qualified types, \ruleref{tc-capt} composes the existing evidence
\(Q\semleq C\) with its premise \(C\semleq D\).  Shape action changes the
shape evidence while retaining the same \(Q\).  This is the step that makes
capture prediction survive subsumption.

\subsection{Action for a dependent pair}
\label{sec:pair-action}

Suppose \(z\) has pair evidence with introduction qualifier \(Q\):
\[
 \Omega;\sigma\models
 z:\capt{\pairshape{x}{T}{a}{\mu}}{C}@Q,
 \qquad
 \Omega;\sigma\vdash Q\semleq C.
\]
Inverting \ruleref{sem-loc-pair} exposes a first location \(y\), a member
referent \(r\), and evidence
\[
 R_y:\Omega;\sigma\models y:T@Q_y,
 \qquad
 R_r:\Omega;\sigma\models r:\subst{\mu}{y}{x}.
\]
The shape component of a compiled pair coercion contains
\[
 C_1:\Omega;\sigma\vdash T\tcoerce T',
 \qquad
 M:\mathsf{Member}_{\Omega;\sigma}(x:T\vdash\mu<:\mu').
\]
Its action first obtains
\(C_1(R_y):\Omega;\sigma\models y:T'@Q_y\).  It then instantiates the saved
source premise in \(M\) with the same location \(y\), using the source
evidence \(R_y\):
\[
 M[y/x](R_y):
 \Omega;\sigma\vdash
 \subst{\mu}{y}{x}\coerce\subst{\mu'}{y}{x}.
\]
Applying this coercion to \(R_r\) realizes the target member.  The outer pair
retains introduction qualifier \(Q\); the capture component of the enclosing
type coercion composes \(Q\semleq C\semleq D\).  This proves the general pair
rule for term, type, and capture members without requiring the dependent
member to ignore \(x\).

Member action creates a coercion only after the outer pair coercion has
started.  Structural recursion on the original coercion is therefore
insufficient.  Pair cells mention only older locations, and stored type and
capture definitions occupy the base stratum.  The recursive member call
moves to a smaller allocation stratum; all other calls decrease coercion
size.  The lexicographic pair of allocation stratum and coercion size is a
well-founded measure.

\section{The qualifier-aware machine invariant}
\label{sec:invariant}

The machine invariant retains result types and use sets together.  Source
subsumption is represented by a type-coercion suffix and a semantic
subcapturing suffix at the outer runtime constructor.  The judgment
\(\Omega;\sigma\vdash_{\mathsf e}t:T\uses C\) assigns the running term its
current type and use bound.  A continuation judgment
\[
 \Omega;\sigma\vdash k:(S,E)\Rightarrow(T,C)
\]
says that a continuation accepting a current result of type \(S\) and use
\(E\) produces final type \(T\) with use bounded by \(C\).

\begingroup
\setlength{\abovedisplayskip}{5pt plus 1pt minus 1pt}
\setlength{\belowdisplayskip}{5pt plus 1pt minus 1pt}
\setlength{\intextsep}{5pt}
\captionsetup{skip=6pt}
\rulesetup

\begin{figure}[H]
\plateheader{Term evidence}
 {\Omega;\sigma\vdash_{\mathsf e}t:T\uses C}
\noindent
\begin{minipage}[t]{0.47\linewidth}
\infrule[\ruledef{te-path}]
  {\sigma\vdash p\resolve\loc x
   \\
   \Omega;\sigma\vdash
     \capt{\sing p}{\cset p}\tcoerce T
   \\
   \Omega;\sigma\vdash\cset p\semleq C}
  {\Omega;\sigma\vdash_{\mathsf e}p:T\uses C}

\infrule[\ruledef{te-value}]
  {\Omega;\sigma\vdash_{\mathsf v}v:T@Q
   \\
   \Omega;\sigma\vdash\cempty\semleq C}
  {\Omega;\sigma\vdash_{\mathsf e}v:T\uses C}
\end{minipage}\hfill
\begin{minipage}[t]{0.51\linewidth}
\infrule[\ruledef{te-app}]
  {\Omega;\sigma\vdash_{\mathsf e}
      p:\capt{(x:S)\to U}{C_f}\uses C_p
   \\
   \Omega;\sigma\vdash_{\mathsf e}q:S\uses C_q
   \\
   \Omega;\sigma\vdash\subst{U}{q}{x}\tcoerce T
   \andalso
   \Omega;\sigma\vdash C_p\cup C_q\semleq C}
  {\Omega;\sigma\vdash_{\mathsf e}p\,q:T\uses C}

\infrule[\ruledef{te-let}]
  {\Omega;\sigma\vdash_{\mathsf e}t:S\uses C
   \\
   \mathsf{Body}_{\Omega;\sigma}
      (x:S\vdash u:U^\uparrow\uses C^\uparrow)
   \\
   \Omega;\sigma\vdash U\tcoerce T
   \andalso
   \Omega;\sigma\vdash C\semleq D}
  {\Omega;\sigma\vdash_{\mathsf e}
    \mathbf{let}\ x=t\ \mathbf{in}\ u:T\uses D}
\end{minipage}

\caption{The complete qualifier-aware runtime term evidence.}
\Description{All rules for runtime terms with result types and use sets.}
\label{fig:term-evidence}
\end{figure}

\begin{figure}[H]
\ContinuedFloat
\plateheader{Continuation evidence}
 {\Omega;\sigma\vdash k:(S,E)\Rightarrow(T,C)}
\noindent
\begin{minipage}[t]{0.39\linewidth}
\infrule[\ruledef{ke-hole}]
  {\Omega;\sigma\vdash S\tcoerce T
   \\
   \Omega;\sigma\vdash E\semleq C}
  {\Omega;\sigma\vdash[\,]:(S,E)\Rightarrow(T,C)}
\end{minipage}\hfill
\begin{minipage}[t]{0.58\linewidth}
\infrule[\ruledef{ke-cons}]
  {\Omega;\sigma\vdash k:(U,D)\Rightarrow(T,C)
   \\
   \mathsf{Body}_{\Omega;\sigma}
      (x:S\vdash u:V^\uparrow\uses F^\uparrow)
   \\
   \Omega;\sigma\vdash V\tcoerce U
   \andalso
   \Omega;\sigma\vdash E\semleq C
   \andalso
   \Omega;\sigma\vdash F\semleq D}
  {\Omega;\sigma\vdash
   (u::k):(S,E)\Rightarrow(T,C)}
\end{minipage}

\plateheader{State evidence}
 {\Omega\models\langle\sigma,k,t\rangle:(T,C)}
\infrule[\ruledef{se-state}]
  {\Omega\valid\sigma
   \andalso
   \Omega;\sigma\vdash k:(S,E)\Rightarrow(T,C)
   \\
   \Omega;\sigma\vdash_{\mathsf e}t:S\uses E}
  {\Omega\models\langle\sigma,k,t\rangle:(T,C)}

\caption{The complete continuation and state evidence.}
\Description{All rules for continuation and machine-state evidence.}
\label{fig:state-evidence}
\end{figure}

\endgroup

Value evidence separates retained captures from immediate uses.  A value has
empty use in \ruleref{te-value}, while its introduction qualifier remains in
the premise.  Application evidence records both operand uses and their common
upper bound.  Continuation evidence carries that bound through suspended let
bodies, so state evidence gives one public use bound for the remainder of the
execution.

\begin{lemma}[Source typing interpretation]
\label{lem:source-interpretation}
If \(\Gamma\vdash t:T\uses C\),
\(\rho\models^\Omega_\sigma\Gamma\), and \(\Omega\valid\sigma\), then
\[
 \Omega;\sigma\vdash_{\mathsf e}
   \rho(t):\rho(T)\uses\rho(C).
\]
\end{lemma}

The proof follows source typing.  Abstraction and let retain their body
derivations, including the use index.  Pair introductions use the
introduction types above.  Applications interpret both paths.  Subsumption
compiles type subtyping and subcapturing and composes both suffixes at the
outer constructor.

\section{Soundness and capture predictions}
\label{sec:soundness}

Write \(s\steps s'\) for the reflexive-transitive closure of machine steps,
and let \(\mathsf{progress}(s)\) mean that \(s\) is final or takes a step.

\begin{theorem}[Semantic progress]
\label{thm:semantic-progress}
If \(\Omega\models s:(T,C)\), then \(\mathsf{progress}(s)\).
\end{theorem}

The term-evidence constructor determines the proof.  Paths resolve, values
return or allocate, and lets push a frame.  In the application case,
qualifier-preserving coercion action exposes a stored abstraction at the
operator location and a well-typed argument at the operand location.

Allocation changes the location scope.  We write \(T\extends U\) and
\(C\extends D\) when the right-hand object is obtained from the left by zero
or more uniform weakenings into freshly extended stores.

\begin{theorem}[One-step preservation]
\label{thm:preservation}
If \(\Omega\models s:(T,C)\) and \(s\sstep s'\), then there are a valid
annotated target store \(\Omega'\), a type \(U\), and a capture set \(D\)
such that
\[
 T\extends U,
 \qquad C\extends D,
 \qquad \Omega'\models s':(U,D).
\]
\end{theorem}

Nonallocating transitions retain \(\Omega,T,C\).  Allocation obtains
\(v:T_v@Q\) from the current value evidence, extends the store annotation by
\(v@Q\), and applies \ruleref{world-allocate}.  All earlier evidence is
weakened once, including the final type and use set, after which the waiting
body is interpreted at the fresh location.

The path case uses co-resolution to convert the source path to the location
returned by \ruleref{e-path}.  The let transitions move a body closure into
the continuation or instantiate it with an existing location.  The
application case is the only case where the introduction qualifier of a
stored value contributes to evaluation; it is detailed next.

\subsection{Beta reduction and use coverage}

Suppose an application \(p\,q\) resolves its operator to a stored closure
with introduction qualifier \(Q\), and resolves its argument to location
\(y\).  The body closure retained by \ruleref{sem-loc-fun} has use
\(Q^\uparrow\cup\cset x\).  Opening the body at \(y\) therefore yields use
\(Q\cup\cset y\).

Term evidence supplies
\[
 \Omega;\sigma\vdash\cset p\semleq C_p,
 \qquad
 \Omega;\sigma\vdash\cset q\semleq C_q,
 \qquad
 \Omega;\sigma\vdash C_p\cup C_q\semleq C.
\]
Since \(p\) names the closure, \ruleref{cap-fold} gives
\(Q\semleq\cset p\).  Since \(q\) and \(y\) resolve to the same location,
\ruleref{cap-alias} gives \(\cset y\semleq\cset q\).  Hence
\[
\begin{array}{rcl}
 Q &\semleq& \cset p\semleq C_p\semleq C,
 \\
 \cset y &\semleq& \cset q\semleq C_q\semleq C,
\end{array}
\]
and \ruleref{cap-union-elim} derives
\(Q\cup\cset y\semleq C\).  The reduced body consequently has the same
public use bound as the application.  The deferred codomain coercion is
instantiated with \(y\), and runtime conversion relates the codomain opened
with \(y\) to the source result opened with \(q\).

\begin{theorem}[Finite preservation]
\label{thm:finite-preservation}
If \(\Omega\models s:(T,C)\) and \(s\steps s'\), then there are
\(\Omega',U,D\) such that \(\Omega'\) is valid for the target store,
\(T\extends U\), \(C\extends D\), and
\(\Omega'\models s':(U,D)\).
\end{theorem}

The empty annotated store validates the empty store.  Source interpretation
with the empty valuation and an identity empty continuation gives initial
state evidence.

\begin{theorem}[Closed progress]
\label{thm:closed-progress}
If \(\ctxempty\vdash t:T\uses C\), then
\(\mathsf{progress}(\mathsf{initial}(t))\).
\end{theorem}

\begin{theorem}[Closed finite preservation]
\label{thm:closed-preservation}
If \(\ctxempty\vdash t:T\uses C\) and
\(\mathsf{initial}(t)\steps s\), then there are
\(\Omega,U,D\) such that \(\Omega\) validates the store in \(s\),
\(T\extends U\), \(C\extends D\), and
\(\Omega\models s:(U,D)\).
\end{theorem}

\begin{theorem}[Type safety]
\label{thm:type-safety}
If \(\ctxempty\vdash t:T\uses C\) and
\(\mathsf{initial}(t)\steps s\), then \(\mathsf{progress}(s)\).
\end{theorem}

This is the ordinary non-stuckness result: every finite execution prefix of a
closed, well-typed program ends in a final state or a state that can step.

\subsection{Use prediction}

An invocation event records the two paths inspected by an application step:
\begingroup
\rulesetup
\infrule[\ruledef{inv-app}]
  {\sigma\vdash p\resolve\loc f
   \andalso
   \sigma\vdash q\resolve\loc y
   \andalso
   \sigma(f)=\lambda(x:A).t}
  {\mathsf{invokes}(
    \langle\sigma,k,p\,q\rangle\sstep
    \langle\sigma,k,\subst{t}{y}{x}\rangle;,p,q)}
\endgroup

\begin{theorem}[Use prediction]
\label{thm:use-prediction}
Suppose \(\ctxempty\vdash t:T\uses C\),
\(\mathsf{initial}(t)\steps s\), and a step from \(s\) invokes \(p\) with
argument path \(q\).  Then there are a valid annotation \(\Omega\) of the
store \(\sigma\) in \(s\) and an allocation extension \(C\extends D\) such
that
\[
 \Omega;\sigma\vdash\cset p\semleq D
 \qquad\text{and}\qquad
 \Omega;\sigma\vdash\cset q\semleq D.
\]
\end{theorem}

The result follows by finite preservation and inversion of
\ruleref{te-app}; continuation evidence lifts the term-local coverage to the
public use bound \(D\).  It concerns precisely the operand paths inspected by
the CK application transition.

\subsection{Capture prediction for returned values}

The final-state relation \(s\returns v\) exposes the returned value whether
it appears directly or through a final location:
\begingroup
\rulesetup
\infrule[\ruledef{returns-value}]
  {v\ \mathsf{value}}
  {\langle\sigma,[\,],v\rangle\returns v}

\infrule[\ruledef{returns-location}]
  {\sigma(x)=v}
  {\langle\sigma,[\,],x\rangle\returns v}
\endgroup

\begin{theorem}[Final capture prediction]
\label{thm:capture-prediction}
Suppose \(\ctxempty\vdash t:T\uses C\),
\(\mathsf{initial}(t)\steps s\), and \(s\) is final.  Then there are a valid
annotation \(\Omega\), allocation extensions \(T\extends U\) and
\(C\extends D\), a returned value \(v\), a type \(V\), and an introduction
qualifier \(Q\) such that
\[
 s\returns v,
 \qquad
 \Omega;\sigma\vdash_{\mathsf v}v:V@Q,
 \qquad
 \Omega;\sigma\vdash Q\semleq\qual(U).
\]
\end{theorem}

For a direct value, the conclusion follows from its value evidence and the
empty continuation's coercion suffix.  For a returned location, validity
supplies the value and introduction qualifier stored at that location;
\ruleref{cap-fold}, the path's type-coercion suffix, and the continuation
suffix relate that qualifier to \(\qual(U)\).

Type safety, use prediction, and final capture prediction have distinct
conclusions.  Type safety rules out stuck states.  Use prediction bounds the
paths inspected by application.  Capture prediction bounds the capabilities
retained by a returned value.  The calculus has no primitive capability
operation or effect event, so these results do not constitute a theorem about
I/O or another external effect semantics.

\section{Limitations and next steps}

The calculus has immutable stores, monadic normal form, stable paths,
dependent functions, singleton types, abstract type and capture members, and
dependent pairs.  It omits intersections, recursive self types, mutable
state, higher-kinded types, and a module initialization discipline.  These
omissions keep the proof independent of the narrowing and inert-type
infrastructure used by larger DOT calculi.

The dynamic observations in this development are application operands and
returned values.  Adding primitive capabilities would require syntax and an
event semantics for their operations, together with a theorem connecting
those events to the use bound.  The metatheoretic coercions would remain proof
objects.

A next extension can add intersections in a controlled fragment and
determine which distributivity and member-merging principles are needed by
Scala encodings.  Recursive objects require a precise account of
initialization and stable self paths.  The general pair proof identifies the
main semantic constraint: coercion action relies on members referring only to
older store locations.  Mutable state or recursive allocation would break
this stratification and require another model, such as a step-indexed
relation.

The mechanization proves bound consistency and its preservation by member
subtyping.  Full context formation and typing regularity remain open because
they require the opening and path-substitution metatheory needed to recover
well-formed synthesized signatures from arbitrary typing derivations.

\section{Related work}

The calculus follows the path-dependent tradition initiated by foundations
for Scala and DOT \citep{amin2014foundations,rompf2016dot}.  Later DOT and
pDOT proofs develop precise typing, inert contexts, and path replacement for
substantially richer calculi
\citep{rapoport2017simple,rapoport2019pdot}.  The restricted type language
here permits a direct store-indexed interpretation.  Step-indexed logical
relations provide another route to full DOT and mutable or recursive
features \citep{giarrusso2020gdot}; the append-only allocation order suffices
for the present system.

Interpreting subtyping as explicit evidence has a long history
\citep{breazutannen1991coercion,crary2000inclusive}.  System FC makes equality
evidence part of an explicitly typed intermediate language
\citep{sulzmann2007systemfc}.  Our coercions remain semantic objects used in
the metatheory, while the source language contains neither coercion terms nor
modalities.  An explicit intermediate elaboration may become useful when
connecting the calculus to a realistic compiler
\citep{martres2023compilation}.

Capture qualifiers and use sets follow Capturing Types
\citep{boruch2023capture}.  The System Capless core includes bounded and
unbounded quantification over capture sets; its surface-language account is
designed to infer this polymorphism and avoid exposing explicit capture
variables in ordinary programs \citep{xu2025box}.  Abstract capture members
provide path-dependent capture-set abstraction here.  Use prediction and
final capture prediction give the corresponding qualifiers an operational
meaning for the observations present in this calculus.

Type and capture members both use intervals.  Their variance follows the
standard order on lower and upper bounds; consistency premises exclude
inverted intervals at formation and selection sites.  More expressive
higher-order interval systems raise additional questions outside the current
calculus \citep{stucki2021intervals}.

\section{Conclusion}

\lambdapcc{} combines path-dependent types and capture checking through one
dependent-pair mechanism.  Abstract type and capture members share path
lookup and interval subtyping; proper term paths additionally support root
contraction.  Store-indexed coercions account for explicit transitivity,
path selections, dependent functions, and the general pair rule.  The
qualifier-aware invariant preserves both result types and use sets while
retaining the introduction qualifier of each stored value.  The mechanized
development proves progress, preservation, and non-stuckness, predicts both
paths inspected by every reachable application, and bounds the introduction
qualifier of every returned value by the qualifier of its result type.

\bibliographystyle{ACM-Reference-Format}
\bibliography{../references}

\end{document}
